% chapter01 —— 由 tools/convert.py 自动生成，勿手改
% 源文件：part0000_split_018.html, part0000_split_019.html, part0000_split_020.html, part0000_split_021.html, part0000_split_022.html

\begin{myEpigraph}{——Graydon Hoare}
在某些情况（例如 Rust 的目标环境）下，比竞争对手快 10 倍，哪怕只快两倍就能成为决胜的关键。速度决定了一个系统在市场上的命运，就像在硬件市场上一样。
\end{myEpigraph}

\begin{myEpigraph}{——《结构化并行程序设计：高效计算模式》，Michael McCool 等}
现在所有的计算机都支持并行……并行编程\textbf{就是}编程。
\end{myEpigraph}

\begin{myEpigraph}{——Andy Wingo}
就连 TrueType 解析器的缺陷都会被攻击者用于监视。安全性对所有软件都很重要。
\end{myEpigraph}

我们用这 3 条引言作为本书的开篇是别有深意的。但还是先从一个“谜题”开始吧。下面的 C 程序是做什么的？

\begin{Verbatim}[commandchars=\\\{\},breaklines=true,breakanywhere=true,fontsize=\codefont,xleftmargin=17bp]
\PYG{n}{int}\PYG{+w}{ }\PYG{n}{main}\PYG{p}{(}\PYG{n}{int}\PYG{+w}{ }\PYG{n}{argc}\PYG{p}{,}\PYG{+w}{ }\PYG{k+kt}{char}\PYG{+w}{ }\PYG{o}{*}\PYG{o}{*}\PYG{n}{argv}\PYG{p}{)}\PYG{+w}{ }\PYG{p}{\PYGZob{}}
\PYG{+w}{  }\PYG{n}{unsigned}\PYG{+w}{ }\PYG{n}{long}\PYG{+w}{ }\PYG{n}{a}\PYG{p}{[}\PYG{l+m+mi}{1}\PYG{p}{]}\PYG{p}{;}
\PYG{+w}{  }\PYG{n}{a}\PYG{p}{[}\PYG{l+m+mi}{3}\PYG{p}{]}\PYG{+w}{ }\PYG{o}{=}\PYG{+w}{ }\PYG{l+m+mh}{0x7ffff7b36ceb}\PYG{n}{UL}\PYG{p}{;}
\PYG{+w}{  }\PYG{k}{return}\PYG{+w}{ }\PYG{l+m+mi}{0}\PYG{p}{;}
\PYG{p}{\PYGZcb{}}
\end{Verbatim}

今天早上，这个程序在 Jim 的笔记本计算机上打印出了下列内容：

\begin{Verbatim}[commandchars=\\\{\},breaklines=true,breakanywhere=true,fontsize=\codefont,xleftmargin=17bp]
\PYG{n}{undef}\PYG{p}{:}\PYG{+w}{ }\PYG{n+nc}{Error}\PYG{p}{:}\PYG{+w}{ }\PYG{p}{.}\PYG{n}{netrc}\PYG{+w}{ }\PYG{n}{file}\PYG{+w}{ }\PYG{n}{is}\PYG{+w}{ }\PYG{n}{readable}\PYG{+w}{ }\PYG{n}{by}\PYG{+w}{ }\PYG{n}{others}\PYG{p}{.}
\PYG{n}{undef}\PYG{p}{:}\PYG{+w}{ }\PYG{n+nc}{Remove}\PYG{+w}{ }\PYG{n}{password}\PYG{+w}{ }\PYG{n}{or}\PYG{+w}{ }\PYG{n}{make}\PYG{+w}{ }\PYG{n}{file}\PYG{+w}{ }\PYG{n}{unreadable}\PYG{+w}{ }\PYG{n}{by}\PYG{+w}{ }\PYG{n}{others}\PYG{p}{.}
\end{Verbatim}

然后它就崩溃了。如果你在自己的机器上尝试运行这个程序，则可能会出现不一样的结果。这是怎么回事呢？

这个程序存在缺陷。数组 \texttt{a} 只有一个元素长，所以根据 C 编程语言标准，使用 \texttt{a[3]} 是\textbf{未定义行为}（undefined behavior）\footnote{对于未定义行为，没有任何标准规定这种情况下该怎么做，因此任何后果都有可能出现。——译者注}：

\begin{myNote}
当使用一个不可移植的或错误的程序结构或者数据时，如果本国际标准对其行为没有强制要求，那么此行为就是未定义的。
\end{myNote}

未定义行为的后果不仅仅是产生难以预料的结果，事实上，它甚至会允许程序做\textbf{任何事}。在我们的例子中，将这个奇怪的数值存储在数组 \texttt{a} 的第四个元素中恰好会破坏 \texttt{main} 函数的调用栈，这样，当从 \texttt{main} 函数返回时，它不会正常退出程序，而是会跳转到 C 标准库中用以从用户主目录的文件中获取密码的代码。这可不妙。

C 和 C++ 给出了数百条规则来规避未定义行为，其中大多数是常识，比如不要访问不应该访问的内存、不要让算术运算溢出、不要除以 0，等等。但是编译器并不会强制执行这些规则——它甚至没有义务去“揭露”这些明显的违规行为。事实上，前面的程序在编译时就没有抛出错误或发出警告。作为程序员，你需要打起十二分的精神来避免这种未定义行为。

从既往经验来看，作为程序员，我们在这方面可没有什么好记录。在犹他大学学习期间，研究员 Peng Li 修改了 C 编译器和 C++ 编译器，让它们“翻译”过的程序能报告自己是否执行了某种形式的未定义行为。他发现几乎所有程序都执行过某种未定义行为，包括那些一直以坚持高标准而声名卓著的项目。幻想靠自己就能避开 C 和 C++ 中的未定义行为，就像幻想只要知道了象棋的规则就能赢得比赛一样。

偶尔出现的奇怪消息或崩溃可能仅仅是质量问题，但自从 1988 年莫里斯（Morris）蠕虫利用类似上述内存越界的技术变体破译用户口令，继而在早期互联网上大规模扩散以来，未定义行为一直是造成各种安全漏洞的主要原因。

因此 C 和 C++ 将程序员置于这样一种尴尬的境地：这些语言是系统编程的行业标准，但它们对程序员的要求实在太高了，这就注定系统会源源不断地出现崩溃和安全问题。要解开这个“谜题”，只会引出一个更大的问题：难道我们不能做得更好吗？

\section{Rust 为你负重前行}

本章开头的 3 条引言给出了我们的答案。第三条引言来自关于震网（Stuxnet）病毒的报道，这是一种计算机蠕虫病毒，2010 年人们发现它入侵了工业控制设备。它使用许多技术来控制受害者的计算机，其中就包括负责解析文字处理软件中内嵌 TrueType 字体的相关代码的未定义行为。可以肯定的是，这段代码的作者并没有料到它以这种方式被利用，这表明需要担心安全问题的不仅仅是操作系统和服务器，任何可能处理“来源不可信数据”的软件都会成为漏洞利用的目标。

Rust 语言给了我们一个简单的承诺：只要程序通过了编译器的检查，就不会存在未定义行为。悬空指针、双重释放和空指针解引用都能在编译期捕获。数组引用则会受到编译期检查和运行期检查的双重保护，因此不会存在缓冲区溢出：想想那个不幸的 C 语言程序，它的 Rust 版代码会安全地退出并给出一条错误消息。

另外，Rust 的目标是既\textbf{安全}又\textbf{易用}。为了更好地保障程序的行为，Rust 对代码施加了比 C 和 C++ 更多的限制，而只有靠实践和经验，我们才能逐渐适应这种限制。但就整体而言，这门语言还是灵活且富有表现力的。用 Rust 编写的代码范围之宽、应用领域之广就是证明。

根据我们的经验，如果该语言能发现更多错误，那么我们就可以尝试一些更具雄心的项目。如果内存管理和指针有效性问题已得到妥善处理，那么修改大型复杂程序时的风险就会降低。如果代码里的 bug 不会破坏程序中的其他部分，那么调试也会简单得多。

当然，仍然有很多 Rust 无法检测到的 bug。但实际上，只要从可能出现的问题清单中排除了未定义行为，就会让编程更美好。

\section{高效并行编程}

众所周知，在 C 和 C++ 中恰当使用并发的难度极大。通常只有在确信单线程代码无法达到预期的性能时，开发人员才会转向并发。但正如本章第二条引言所指出的，并行性对现代机器来说太重要了，不应该等到迫不得已时才考虑它。

事实证明，Rust 用来确保内存安全的那些限制同样能确保 Rust 程序避免产生数据竞争（data race）。只要数据不可变，你就可以在线程之间自由地共享这些数据。会发生变化的数据则只能使用同步原语访问。所有传统的并发工具仍然可用：互斥锁、条件变量、通道、原子等。Rust 只负责检查你是否正确地使用了它们。

这就让 Rust 成了一门能充分发挥现代多核机器能力的优秀语言。Rust 的生态系统提供了一些超乎于常规并发原语的库，可帮助你在处理器池之间均匀分布复杂负载、使用无锁同步机制（如读取-复制-更新）等。

\section{性能毫不妥协}

最后，正如本章第一条引言所说的那样，Rust 同样怀揣着 Bjarne Stroustrup 在他的论文“Abstraction and the C++ Machine Model”（抽象和 C++ 机器模型）中为 C++ 阐明的雄心。

\begin{myNote}
通常，C++ 的各种实现会遵循零开销原则：没用到的，就没有开销；要用到的，你也无法手写出更好的代码。
\end{myNote}

系统编程通常会涉及极限压榨机器性能。对电子游戏来说，整台机器应该极力为玩家创造最佳体验。对浏览器来说，浏览器的效率决定了内容作者可用能力的上限。在机器的固有限制下，必须将尽可能多的内存和处理器能力留给内容本身。同样的原则也适用于操作系统：内核应该尽可能把机器资源留给用户程序，而不是自己来消耗它们。

但是当我们说 Rust“快”的时候，到底意味着什么呢？毕竟人们可以用任何通用语言来写出慢速代码。更准确一点儿说，如果你准备让程序充分发挥底层机器的能力，那么 Rust 就会为你提供支持。这门语言设计了一些“高性能”的默认选项，并赋予你自主控制内存使用和处理器算力分配方式的能力。

\section{协作无边无界}

本章的标题中还隐藏了第四条引言：“系统程序员也能享受美好”。这是指 Rust 对代码共享和复用的支持。

作为 Rust 的包管理器和构建工具，Cargo 能让你轻松使用别人在 Rust 的公共包存储库 crates.io 网站上发布的各种库。只需将库的名称和所需的版本号添加到文件中，Cargo 就会负责下载这个库以及它所用到的任何其他库，并将所有内容链接在一起。你可以将 Cargo 视为 Rust 下的 NPM 或 RubyGems，其侧重于实现健全的版本管理和可重现的构建。一些广为流行的 Rust 库提供了包括开箱即用的序列化功能、HTTP 客户端和服务器功能以及现代图形 API 功能在内的一切。

再进一步说，Rust 这门语言本身也旨在支持协作：借助 Rust 的特型（trait）\footnote{\textbf{特型}，表示一组具有指定特点的类型。这里之所以把 trait 译为“特型”，是因为特型和泛型是紧密相关的，我们希望读者顾名思义就能在两者之间建立联系；而其他中文译法，或者缺乏特异性，或者难以望文生义，因此并未采纳。关于 trait 翻译与否及译法是个非常有争议的话题，我们在审读群曾经进行过至少 3 次深度讨论，但并未形成共识。对于一个常用专有名词，有一个合适的中文表述是非常必要的。但对于一个词的新译法，一开始，大家可能还不太习惯，不妨试着给新词些时间。——译者注}和泛型，我们可以创建具有灵活接口的库，将其用在许多不同的上下文中。Rust 的标准库提供了一组最核心的基本类型，这些类型为一些常见的情况建立了共享规约，以方便不同的库彼此协作。

第 2 章旨在通过介绍几个小的 Rust 程序来具体展示本章所述的这几点优势。

